% Part: history
% Chapter: biographies
% Section: bertrand-russell

\documentclass[../../../include/open-logic-section]{subfiles}

\begin{document}

\olfileid{his}{bio}{rus}

\olsection{برتراند راسل}

\olphoto{russell-bertrand}{Bertrand Russell}

يُحتفى ببرتراند راسل بوصفه أحد مؤسسي الفلسفة التحليلية الحديثة. وُلد في
18 مايو 1872، ولم يُعرف بعمله في الفلسفة والمنطق فحسب، بل كتب أيضاً
مؤلفات شعبية كثيرة في موضوعات متنوعة. وكان كذلك ناشطاً سياسياً متحمساً
طوال حياته.

وُلد راسل في تريليك بمونماوثشير في ويلز. وكان والداه من النبلاء
البريطانيين. وكانا من دعاة الفكر الحر، بل صادقا المتطرفين في بوسطن آنذاك.
لكن والدي راسل توفيا، للأسف، وهو صغير، فأُرسل ليعيش مع جديه. وهناك تلقى
تنشئةً دينية (وهو ما أراد والداه تجنبه بأي ثمن). وكانت جدته شديدة الصرامة
في جميع مسائل الأخلاق. وفي مراهقته تلقى معظم تعليمه في المنزل على أيدي
معلمين خصوصيين.

كان تأثير راسل في الفلسفة التحليلية، ولا سيما المنطق، هائلاً. درس
الرياضيات والفلسفة في كلية ترينيتي بجامعة كامبريدج، وتأثر هناك بالرياضي
والفيلسوف ألفرد نورث وايتهد. وفي عام 1910 نشر راسل ووايتهد المجلد الأول
من \emph{مبادئ الرياضيات}، ودافعا فيه عن الرأي القائل إن الرياضيات قابلة
للرد إلى المنطق. ثم نشر مئات الكتب والمقالات والمنشورات السياسية. وفي
عام 1950 نال جائزة نوبل في الأدب.

كان راسل منغمساً بعمق في السياسة والنشاط الاجتماعي. فخلال الحرب العالمية
الأولى اعتُقل وسُجن ستة أشهر بسبب أنشطته السلمية واحتجاجاته. واستطاع في
السجن أن يكتب ويقرأ، وزعم أنه وجد التجربة «مقبولةً جداً». وظل سلمياً طوال
حياته، وسُجن مرة أخرى لحضوره تجمعاً لنزع السلاح النووي عام 1961. كما نجا
من تحطم طائرة عام 1948، ولم ينجُ فيه سوى الجالسين في قسم المدخنين. ولذلك
ادعى راسل أنه مدين بحياته للتدخين. تزوج أربع مرات، وكانت له سمعة في إقامة
علاقات خارج الزواج. وتوفي في 2 فبراير 1970 عن عمر 97 عاماً في
بنرينديودرايث بويلز.

\begin{reading}
كتب راسل سيرةً ذاتية في ثلاثة أجزاء تغطي حياته من 1872 إلى 1967
\citep{Russell1967,Russell1968,Russell1969}. ويضم مركز برتراند راسل
للبحوث في جامعة ماكماستر أرشيف برتراند راسل. انظر موقعه عند
\citet{Duncan2015} للحصول على معلومات عن مجلدات أعماله المجموعة (ومنها
كشافات قابلة للبحث) وعن مشروعات الأرشفة. ويعد بحث راسل
\emph{في الدلالة} \citep{Russell1905} من كلاسيكيات الفلسفة التحليلية في
القرن العشرين.

يضم مدخل موسوعة ستانفورد للفلسفة عن راسل \citep{Irvine2015} مقاطع صوتية
له يتحدث فيها عن الرغبة والنظرية السياسية. وتتوفر على الإنترنت مقابلات
مصورة كثيرة معه. ولرؤيته يتحدث، مثلاً، عن التدخين ووقوعه في حادث تحطم
طائرة، انظر \citet{RussellND}. وتتوفر بعض أعمال راسل، ومنها
\emph{مقدمة في الفلسفة الرياضية}، كتباً صوتية مجانية لدى
\citet{LibriVoxND}.
\end{reading}

\end{document}
